% !TeX root = ../papers/nc-p3.tex
\section{対称完全交叉の全点での三階数切断の滑らかさ}
\label{sec:ci-global-tangent}
$\mathfrak X_6$は三つの二対二分割について全ての7次小行列式を課した射影スキームである。
定理\ref{thm:symmetric-presentation-product}の対称チャート$S$が、そのどの点でも
$\mathfrak X_6$の局所構造全体を与えることを示す。

\begin{lemma}[中央三次数の次数零剛性]\label{lem:ci-middle-rigidity}
$B=\operatorname{Sym}(V^*)/(L)$を四つの二次式のArtin完全交叉とし、
$\mathfrak m=B_{>0}$, $N=\mathfrak m/B_4$を元の次数$1,2,3$に置く。
このとき$\operatorname{End}_B(N)_0=k$、
$\operatorname{Ext}^1_B(N,N)_0=0$である。
\end{lemma}
\begin{proof}
全てのHomとExtで内部次数零を取る。$D_0=B/B_4$と置く。
$B$はFrobeniusで自己入射的、$B_4=k(-4)$であり、
二次完全交叉はKoszulなので
$\operatorname{Ext}^p_B(k,k(-4))_0=0$ ($p\ne4$)である。
$0\to k(-4)\to B\to D_0\to0$から
$\operatorname{Ext}^p_B(k,D_0)_0=0$ ($p=1,2$)を得る。
また$\operatorname{Hom}_B(k,D_0)_0=0$である。
$0\to N\to D_0\to k\to0$と$\operatorname{Ext}^p_B(k,k)_0=0$ ($p>0$)より
$\operatorname{Ext}^1_B(k,N)_0=k$、$\operatorname{Ext}^2_B(k,N)_0=0$となる。
次に$0\to\mathfrak m\to B\to k\to0$で$N_0=0$を使うと
$\operatorname{Hom}_B(\mathfrak m,N)_0=k$、
$\operatorname{Ext}^1_B(\mathfrak m,N)_0=0$である。
最後に$0\to k(-4)\to\mathfrak m\to N\to0$へ$\operatorname{Hom}_B(-,N)_0$を適用する。
$N_4=0$なので$\operatorname{Hom}_B(k(-4),N)_0=0$となり、両方の主張が従う。
\end{proof}

\begin{theorem}[対称チャートは三階数切断の開部分]\label{thm:symmetric-open-all}
$\mathfrak X_6$は$S$の全点で滑らかであり、$S$は$\mathfrak X_6$の69次元の開部分である。
特に$Z_{\mathrm{sym}}$と異なる既約成分の交わりは$S$を避ける。
\end{theorem}
\begin{proof}
四因子の基底変更により$[w]=[w_L]$を完全対称なCIのtraceとしてよい。
アフィン錐の非零点$w$での接空間をそれぞれ$T_X,T_S$とする。
まず$\dot w\in T_X$が因子$0,1$の交換で不変とする。
$D=k[\epsilon]/(\epsilon^2)$上の$w'=w+\epsilon\dot w$はこの二因子に対称である。
各分割で階数6の小行列式の一つが単元で、全7次小行列式が消えるため、
行列の像・核は直和因子である（単元小行列で掃き出す）。
$01|23$の縮約の$\operatorname{Sym}^2 V^*$内の核は$L$を持ち上げる階数4の$L'$となる。
$B'=\operatorname{Sym}_D(V^*\otimes D)/(L')$は完全交叉族の引戻しであり、
各次数は$D$上自由である。これは、特殊ファイバーで正則列を解消するKoszul複体を
次数ごとに単元掃き出しすることからも従う。

$02|13$の階数分解を$P'(x)Q'(y)$と書く。第0・第1因子の対称性より
$P'(x)Q'(y)=P'(y)Q'(x)$である。従って順に$Q',P'$を一次式の作用とする
階数$(4,6,4)$の三層加群を作れる。$L'$の全二次式が両作用の積を消し、
その加群は$B'$上である。特殊ファイバーは$N$に同型である。
実際、標準の階数分解は$B_1\to B_2\to B_3$の二つの乗法であり、
$B_3$を$B_1$の双対と同定するtraceは最上テンソルを与える。
全階数分解は中間空間の基底変更でこれに移る。

同じ$B'$に対して標準加群$N'_{\rm st}=B'_{>0}/B'_4$がある。
三次数ごとに基底を合わせ、二つの加群の特殊ファイバーを同一にする。
その二つの一次作用の差は$\epsilon$の係数として$B$上の次数付き自己拡大を定める。
具体的には、$B'$の元$b'$の二つの作用の差を$\epsilon$で割れば、
特殊化$b$だけに依存する写像$\delta(b)$が得られる。
$\delta(bc)=b\delta(c)+\delta(b)c$なので、対角に元の作用、
上三角成分に$\delta$を入れた$N\oplus N$が拡大加群である。
補題\ref{lem:ci-middle-rigidity}から次数零で分裂し、差は次数別基底変更
$1+\epsilon f$で消える。従って$w'$は$B'$の対称traceテンソルから、
入力・出力の基底変更とtraceの単元倍によって得られる。中間次数の基底変更は
$P'Q'$の積で相殺する。ゆえに$\dot w\in T_S$である。

$w$は完全対称、両方の錐は四因子の置換で不変なので$T_X,T_S$は
$\mathfrak S_4$表現である。上の議論を共役すれば、各互換の固定部分は$T_S$に入る。
従って商$T_X/T_S$では全互換が$-1$で作用し、交代化射影が恒等になる。
一方、$T_X$に完全交代テンソルは存在しない。実際、$01|23$行列の左右核は
ともに$\bigwedge^2 V^*$を含む。接条件は左右核上で$\dot w$の双線型形式が零となることである。
完全交代テンソルを二組の交代二形式に評価すると、その四文字係数の4倍を得るので、
この条件は全係数を消す。従って交代化射影は$T_X$上零であり、$T_X=T_S$である。

射影化後の接次元は定理\ref{thm:symmetric-presentation-product}より69である。
$S$を含むので局所次元は少なくとも69であり、接次元との一致から
$\mathfrak X_6$はここで正則かつ滑らかである。通過する既約成分は
$Z_{\mathrm{sym}}$一つで、$S$はその中で開である。他の有限個の成分を除けば
各点の近傍を$S$内に取れるため、全体としても$S$は開である。
\end{proof}

この定理は全ての対称正則点を扱うが、決定式条件を外したAS表示空間の滑らかさや、
$\mathfrak X_6$の他の成分の不存在は与えない。
$S$の閉包には既知の有限Frobenius非Koszul境界がある。

\begin{corollary}[未分類低階数部分の特殊化]\label{cor:closed-low-residual}
有限Frobenius標準候補のうち、$\rho\leq6$で、対称CIと可換多項式軌道のいずれにも
入らない部分は、全有限Frobenius標準候補の中で相対閉である。
標準AS対象に制限しても同じ相対閉性が成り立つ。
\end{corollary}
\begin{proof}
二つの既知部分は、定理\ref{thm:symmetric-open-all}と
定理\ref{thm:pure-anti-and-alternating-local}により$\mathfrak X_6$の開部分である。
その補集合は$\mathfrak X_6$で閉であり、$\mathfrak X_6$自体が射影テンソル空間で閉なので、
候補部分・AS部分との交叉を取ればよい。未分類AS部分の非空性はこの議論からは従わない。
\end{proof}

\paragraph{平方を含まない完全交叉での検算。}
巡回添字で$f_i=x_i^2+c_ix_{i+1}x_{i+2}$、$(c_0,c_1,c_2,c_3)=(1,2,3,5)$と取る。
共通零点に一つでも非零座標があれば全座標が非零になり、四式の積から$\prod c_i=1$が必要になる。
また$\ell^2\in L$なら、非零座標$\ell_i$は混合項によって$\ell_{i+2}\ne0$を強制するが、
$x_ix_{i+2}$はどの$f_j$にも現れない。従って射影基点も非零平方もない。
\path{computations/verify_ci_middle_rigidity.py}で四次traceを求め、
固定した$B$の二作用を線型化すると192変数で階数125を得る。
核次元67は三次数の基底変更群の次元からスカラー安定化群を除いた数に等しい。
三分割全部のアフィン接次元は70、その$01$交換固定部分は55で、射影接次元は69である。
